\section{Discussion and Conclusion}
\subsection{What the organization retains}
An EEO retains procedures and resource relationships that can be applied beyond the execution that motivated them. This differs from keeping a conversation transcript available to the next task: the revision changes the instructions, examples, or domain references that define how the organization works. The transcript supplies evidence for a change; the package is the object that carries it forward. Its exact identity makes that distinction observable. A successful one-off repair with no retained package change does not instantiate this update mechanism.

The package boundary also exposes a substantive design choice. Procedures, role responsibilities, capability references, and supporting files are revised as one attributable object. This makes it possible to ask whether an improvement depends on a particular professional configuration, rather than only whether a new prompt scores better. Packaging alone does not produce that improvement. The matched GEPA control tests whether the particular revision policy adds value on the same substrate. Comparing the fixed generic team with the fixed specialized team isolates specialization; comparing the fixed specialized team with EEO--text isolates retained updating. A direct comparison between the generic team and EEO--text combines both changes.

\subsection{When specialization helps or hinders}
Specialization is most plausibly useful when a task family shares recurring operations and acceptance criteria: source assessment, population reconciliation, computation, revision-specific verification, and evidence-preserving rendering. The package can place these obligations at the roles and handoffs where they matter. The same structure can be costly when the task requires only faithful rewriting of supplied material. Research Studio's alternative workflow declarations make this variation explicit, but the orchestrator still has to choose an appropriate entry point. Organizational sophistication is therefore not measured by the number of experts invoked.

A retained procedure can improve one failure mode while degrading another. Additional verification may reduce unsupported quantitative claims while increasing delay or abstention. A domain-specific reference can improve familiar tasks and introduce irrelevant assumptions into a new family. Retention evaluation must consequently cover previously useful task types, with the same outcome definitions and recorded resource costs. A mean gain across tasks does not establish that every task or every professional capability improved.

\subsection{Limits of the update boundary}
Frozen executables and inherited capability references make candidate differences inspectable, but can exclude the tool or algorithm change needed to resolve a failure. Such a failure calls for a separately authorized change to the professional substrate, followed by a new comparison context. It cannot be interpreted as evidence that further text search must eventually solve the problem. Within the allowed space, a text edit can still substantially alter behavior or misuse an existing capability; structural validity does not establish semantic safety.

The evaluation boundary imposes another practical cost. Once outcomes influence candidate selection, they become development evidence. Reusing the same test tasks across successive organizational revisions can make repeated apparent improvement an artifact of adaptation to the test set. Fresh partitions and independent search episodes are needed for the corresponding generalization and method-stability claims. The resource cost of collecting and evaluating that evidence belongs alongside the computational cost of search.

\subsection{Conclusion}
EEO makes reusable professional expertise an explicit organization of roles, procedures, capabilities, and handoffs. Domain specialization establishes its working resources and acceptance contracts; experience-driven revision changes how subsequent tasks are performed. OpenCorvus supplies the package, execution-binding, candidate-validation, and comparative-evidence interfaces that connect these changes. The analysis identifies the full-cost and independent-evaluation conditions under which a retained revision can be assessed as useful. The resulting experimental design separates expert collaboration, professional specialization, and retained updating, and evaluates their effects on future-task quality, resource use, and prior capabilities.
